% SPDX-FileCopyrightText: Copyright (c) 2025-2026 Objectionary.com
% SPDX-License-Identifier: MIT

To reason about and transform bytecode structurally, we encode it as \(\varphi\)-expressions,
  a minimal, uniform representation that exposes both instructions and metadata in a single composable form.
The full calculus was introduced by \citet{bugayenko2021eolang}; we use a subset sufficient to represent Java bytecode.

\Cref{fig:jeo} shows a stream pipeline compiled to bytecode and then disassembled into a \(\varphi\)-expression.
Because \phic{} uses composable symbolic structures, its expressions are amenable to rule-based rewriting.
For example, the following rule transforms \ff{IMUL} to \ff{IADD}, where \(B_1\) and \(B_2\) match any sequence of key–value pairs:
\begin{phiquation*}
 \frg{pattern}{ [[ B_1, @ -> Q.imul, B_2 ]] } \longrightarrow \frg{replacement}{ [[ B_1, @ -> Q.iadd, B_2 ]] }.
\end{phiquation*}
Rules may use functions from \phino{}, our rewriting engine.
The following rule renames classes matching a pattern:
\begin{phiquation*}
 [[ B_1, name -> e_1, B_2 ]] \longrightarrow [[ B_1, name -> e_2, B_2 ]]
   \qquad\qquad\text{when}\quad |matches| {(} |"\char94Nums\textdollar{}"| {,} e_1 {)}
   \qquad\qquad\text{where}\quad e_2 := |concat| {(} |"Op"| {,} |name| {)} {.}
\end{phiquation*}
Rules are applied repeatedly until none match; order is irrelevant.

We scan bytecode files using the ASM library\footnote{\url{https://asm.ow2.io/}} and disassemble them into \(\varphi\)-expressions (one \ff{.class} file produces one \ff{.phi} file).

\begin{figure*}
\setthreecapwidths{0.31}{0.27}
\begin{subfigure}[t]{\firstsf}
\begin{ffcode}
import
  java.util.stream.*;
class Numbers {
 int total(int[] a) {
  return IntStream
   .of(a)
   .map((*@\highlight{x -> x * x}@*))
   .sum();
 }
}
\end{ffcode}
\end{subfigure}
\hfill
\begin{subfigure}[t]{\secondsf}
\begin{ffcode}
public Numbers():
 aload_0
 invokespecial
 return
public int
  total(int[]):
 aload_1
 invokestatic
 invokedynamic
 invokeinterface
 invokeinterface
 ireturn
private static int
  (*@\highlight{total\textdollar{}0}@*)(int):
 iload_0
 iload_0
 imul
 ireturn
\end{ffcode}
\end{subfigure}
\hfill
\begin{subfigure}[t]{\thirdsf}
\begin{phiquation*}
[[ |j\textdollar{}Numbers| -> [[
  @ -> Q.class,
  version -> 61,
  access -> 33,
  supername ->
    "java/lang/Object",
  interfaces -> [[ ]],
  |j\textdollar{}object\char64init\char64| -> [[\dots]],
  |j\textdollar{}total| -> [[\dots]],
  |\highlight{j\textdollar{}total\textdollar{}0}| -> [[
    @ -> Q.method,
    access -> 4106,
    descriptor -> "(I)I",
    body -> [[
      @ -> Q.seq.of6,
      i2 -> [[ @ -> Q.iload, |v| -> 0 ]],
      i3 -> [[ @ -> Q.iload, |v| -> 0 ]],
      i4 -> [[ @ -> Q.imul ]],
      i5 -> [[ @ -> Q.ireturn ]] ]] ]] ]] ]]
\end{phiquation*}
\end{subfigure}
\threefigcaps
  {jeo}
  {Java class.}
  {Bytecode, from \ff{javap}.}
  {Simplified \(\varphi\)-expression.}
  {Java class after compilation and disassembly to \(\varphi\)-expression.}
\end{figure*}

Each instruction becomes a formation qualified by a dispatch (e.g., $Q.class$ for classes, $Q.method$ for methods); parameterized instructions like \ff{ILOAD} carry their parameters as attributes.
Java names are prefixed with \ff{j\textdollar{}} to distinguish them from meta-attributes.

The Java compiler translates lambda expressions into \ff{invokedynamic} and \ff{invokeinterface} instruction pairs~\citep{rose2009bytecodes}.
We re\-write these pairs into \pragma{filter} or \pragma{map} pragmas---\phic{} formations that temporarily carry rewriting data but don't exist in the final bytecode.

Both \pragma{filter} and \pragma{map} are then unified into \pragma{distill}, an abstraction accepting an item and returning \ff{void}.
Consecutive \pragma{distill} pragmas are merged by concatenating their lambda bodies.
Finally, \pragma{distill} pragmas become \pragma{mapMulti} pragmas, which generate new private static methods containing the merged logic plus a call to the \ff{Consumer}'s \ff{accept} method.
The pragmas are then transformed back into instructions and reassembled into bytecode via ASM.
